% GENERATED FILE. Edit canonical lessons, then run npm run generate:books.
% canonical-source-hash: fnv1a64:2515c69ab5978908
% canonical-lessons: RU-C154-poetomu, RU-C154-tak-kak, RU-C154-yesli, RU-C154-khotya, RU-C154-inache

\chapter{So, Since, If, Although, Otherwise}
\label{ch:ru-so-since-if-although-otherwise}
\hlchaptermodality{154}

\begin{chapteropening}
\textbf{By the end of this chapter:} \emph{I can say so, since, if, although and otherwise.}

Complete the last lesson of chapter 154: I can say so, since, if, although and otherwise.
\end{chapteropening}

\section[\texorpdfstring{\ru{поэтому} (poétomu)}{поэтому (poétomu)}]{\ru{поэтому} (poétomu) — so, therefore}

\label{lesson:RU-C154-poetomu}



\begin{quote}
Before the new one: say the Russian for a monument, then the Russian for a fountain.
\end{quote}

\subsection*{You'll want to know: \ru{поэтому}}
\textbf{\ru{поэтому}} (\emph{poétomu}) — \textquotedblleft{}so, therefore\textquotedblright{}.

\begin{grammarlens}[title={What you've built}]
One more word for this chapter.
\end{grammarlens}

\subsection*{Your turn}
\begin{itemize}
\raggedright
  \item \emph{Say it:} \emph{poétomu}
  \item \emph{Say it:} \emph{poétomu}, once more
  \item \emph{Say it:} say \emph{poétomu}
  \item \emph{Recall:} say the Russian for a monument, then the Russian for a fountain, then say \emph{poétomu} again
\end{itemize}

\begin{tcolorbox}[breakable,colback=teal!4,colframe=teal!35!black,title={Before you move on}]
What does \ru{поэтому} mean? (\textquotedblleft{}So, therefore\textquotedblright{}.) Say it once more.
\end{tcolorbox}
\section[\texorpdfstring{\ru{так} \ru{как} (tak kak)}{так как (tak kak)}]{\ru{так} \ru{как} (tak kak) — since, as}

\label{lesson:RU-C154-tak-kak}



\begin{quote}
Before the new one: say the Russian for a fountain, then the Russian for so.
\end{quote}

\subsection*{You'll want to know: \ru{так} \ru{как}}
\textbf{\ru{так} \ru{как}} (\emph{tak kak}) — \textquotedblleft{}since, as\textquotedblright{}.

\begin{grammarlens}[title={What you've built}]
One more word for this chapter.
\end{grammarlens}

\subsection*{Your turn}
\begin{itemize}
\raggedright
  \item \emph{Say it:} \emph{tak kak}
  \item \emph{Say it:} \emph{tak kak}, once more
  \item \emph{Say it:} say \emph{tak kak}
  \item \emph{Recall:} say the Russian for a fountain, then the Russian for so, then say \emph{tak kak} again
\end{itemize}

\begin{tcolorbox}[breakable,colback=teal!4,colframe=teal!35!black,title={Before you move on}]
What does \ru{так} \ru{как} mean? (\textquotedblleft{}Since, as\textquotedblright{}.) Say it once more.
\end{tcolorbox}
\section[\texorpdfstring{\ru{если} (yésli)}{если (yésli)}]{\ru{если} (yésli) — if}

\label{lesson:RU-C154-yesli}



\begin{quote}
Before the new one: say the Russian for so, then the Russian for since.
\end{quote}

\subsection*{You'll want to know: \ru{если}}
\textbf{\ru{если}} (\emph{yésli}) — \textquotedblleft{}if\textquotedblright{}.

\begin{grammarlens}[title={What you've built}]
One more word for this chapter.
\end{grammarlens}

\subsection*{Your turn}
\begin{itemize}
\raggedright
  \item \emph{Say it:} \emph{yésli}
  \item \emph{Say it:} \emph{yésli}, once more
  \item \emph{Say it:} say \emph{yésli}
  \item \emph{Recall:} say the Russian for so, then the Russian for since, then say \emph{yésli} again
\end{itemize}

\begin{tcolorbox}[breakable,colback=teal!4,colframe=teal!35!black,title={Before you move on}]
What does \ru{если} mean? (\textquotedblleft{}If\textquotedblright{}.) Say it once more.
\end{tcolorbox}
\section[\texorpdfstring{\ru{хотя} (khotyá)}{хотя (khotyá)}]{\ru{хотя} (khotyá) — although}

\label{lesson:RU-C154-khotya}



\begin{quote}
Before the new one: say the Russian for since, then the Russian for if.
\end{quote}

\subsection*{You'll want to know: \ru{хотя}}
\textbf{\ru{хотя}} (\emph{khotyá}) — \textquotedblleft{}although\textquotedblright{}.

\begin{grammarlens}[title={What you've built}]
One more word for this chapter.
\end{grammarlens}

\subsection*{Your turn}
\begin{itemize}
\raggedright
  \item \emph{Say it:} \emph{khotyá}
  \item \emph{Say it:} \emph{khotyá}, once more
  \item \emph{Say it:} say \emph{khotyá}
  \item \emph{Recall:} say the Russian for since, then the Russian for if, then say \emph{khotyá} again
\end{itemize}

\begin{tcolorbox}[breakable,colback=teal!4,colframe=teal!35!black,title={Before you move on}]
What does \ru{хотя} mean? (\textquotedblleft{}Although\textquotedblright{}.) Say it once more.
\end{tcolorbox}
\section[\texorpdfstring{\ru{иначе} (ináche)}{иначе (ináche)}]{\ru{иначе} (ináche) — otherwise}

\label{lesson:RU-C154-inache}



\begin{quote}
Before the new one: say the Russian for if, then the Russian for although.
\end{quote}

\subsection*{You'll want to know: \ru{иначе}}
\textbf{\ru{иначе}} (\emph{ináche}) — \textquotedblleft{}otherwise\textquotedblright{}.

\begin{grammarlens}[title={What you've built}]
One more word for this chapter.
\end{grammarlens}

\subsection*{Your turn}
\begin{itemize}
\raggedright
  \item \emph{Say it:} \emph{ináche}
  \item \emph{Say it:} \emph{ináche}, once more
  \item \emph{Say it:} say \emph{ináche}
  \item \emph{Recall:} say the Russian for if, then the Russian for although, then say \emph{ináche} again
\end{itemize}

\begin{tcolorbox}[breakable,colback=teal!4,colframe=teal!35!black,title={Before you move on}]
What does \ru{иначе} mean? (\textquotedblleft{}Otherwise\textquotedblright{}.) Say it once more.
\end{tcolorbox}
